\chapter*{Resumen}
\addcontentsline{toc}{chapter}{Resumen}

\begin{summarybox}
Se implementó en CoppeliaSim y Lua un Pioneer P3DX dentro de
una celda de almacén. La escena base \path{Sim_T2_Phase1_Basic.ttt} integra los
elementos principales del escenario: R1, B1 (Bill) como persona móvil, T1 y T2,
dos mesas, estanterías, un sofá de operador, obstáculos y la estación de carga. B1
sale en el código como \texttt{/B1} o Bill; viene del modelo de persona de
CoppeliaSim y se usa como \texttt{People/mannequin}.
\end{summarybox}

R1 lee los sensores de proximidad, estima su pose, actualiza un mapa local,
evita obstáculos y manda las ruedas diferenciales mientras coordina entregas con
B1. El ciclo que se validó fue este.

\begin{itemize}
  \item \texttt{task1:\{T1 -> B1@WS1\}}: R1 recoge T1 en la estantería y la entrega a B1.
  \item \texttt{task2:\{B1@WS1(T1) -> Rack\_T1\}}: B1 trabaja la pieza y R1 la devuelve a su estantería.
  \item \texttt{task3:\{T2 -> B1@WS2\}}: B1 camina a la segunda mesa y solicita T2.
  \item \texttt{task4:\{B1@WS2(T2) -> Rack\_T2\}}: R1 recoge la pieza terminada y la almacena.
\end{itemize}

La segunda escena validada es \path{Sim_T2_Phase2_SLAM_Unknown.ttt}. En esa escena, el
robot empieza con parte del entorno oculta y se encuentra obstáculos que aparecen
durante la misión. En esa fase se comparan cuatro modos de estimación y mapeo escritos en
Lua:
\texttt{KALMAN\_LANDMARK}, \texttt{GMAPPING\_GRID},
\texttt{HECTOR\_GRID\_MATCHING} y \texttt{CARTOGRAPHER\_SUBMAP}. Usan las
mismas tareas, sensores y controlador, para que la comparación no dependa de
condiciones de escena distintas.

La ejecución final de Phase 2 completó las cuatro tareas, volvió a carga, mantuvo
activos los 16 sensores y cerró con \SI{0.040}{m} de error final de pose. En el
ensayo con mapa desconocido, Kalman-landmark quedó con el menor RMSE de posición
(\SI{0.03185}{m}) y la menor duración (\SI{124.80}{s}). Cartographer-submap
generó más actividad de mapa, con 13 submapas y 4 cierres de ciclo locales;
Hector terminó con el mejor indicador aproximado de precisión entre las grids.
